\documentclass[11pt,a4paper]{article}
\usepackage[utf8]{inputenc}
\usepackage[T1]{fontenc}
\usepackage[brazil]{babel}
\usepackage[margin=2.5cm]{geometry}
\usepackage{mathptmx}
\usepackage{amssymb}
\usepackage{enumitem}
\usepackage{booktabs}
\usepackage{tabularx}
\usepackage{array}
\usepackage{forest}
\usepackage{hyperref}
\hypersetup{colorlinks=false, pdfborder={0 0 0}}
\usepackage{fancyhdr}

\pagestyle{fancy}
\fancyhf{}
\fancyhead[L]{\small TrackFleet --- Escopo}
\fancyhead[R]{\small\thepage}
\renewcommand{\headrulewidth}{0.4pt}

\newcolumntype{L}{>{\raggedright\arraybackslash}X}
\newcommand{\dodbox}{$\square$}

\setlist{topsep=3pt}

\begin{document}

\begin{center}
{\Large\bfseries DOCUMENTO DE ESCOPO DO PROJETO}\\[2pt]
{\large\bfseries TrackFleet --- Gestão de Rastreador de Automóveis}\\[4pt]
Disciplina de Gerenciamento de Projetos $\cdot$ Framework PMBOK\textregistered\ 8\textsuperscript{a} Edição $\cdot$ Domínio: Escopo\\[2pt]
Integrantes: Enzo Allebrand e Kerlon Ribeiro (dois GPs) $\cdot$ 4\textsuperscript{o} semestre $\cdot$ Data: 20/08
\end{center}

\vspace{2mm}
\hrule
\vspace{4mm}

\noindent No PMBOK\textregistered\ 8, escopo não é uma coisa só: são três dimensões que acontecem ao mesmo tempo. Este documento reúne as três em um só lugar, cada uma em sua parte, para evitar a armadilha de comprimi-las em um único bloco:
\begin{itemize}[itemsep=2pt]
  \item \textbf{Escopo do Produto (o quê):} as características, funções e atributos do que o TrackFleet deve entregar.
  \item \textbf{Escopo do Projeto (o como):} todo o trabalho necessário, e apenas o necessário, para entregar o produto.
  \item \textbf{Qualidade:} a terceira dimensão do escopo --- não um domínio à parte, mas um atributo integrante dele.
\end{itemize}

\vspace{2mm}
\hrule
\vspace{4mm}

%======================================================================
\section{Escopo do Produto (o quê)}
%======================================================================

\noindent\textbf{\large Escopo do Produto}

\vspace{1mm}
\noindent A descrição das características, funções e atributos do resultado que o TrackFleet deve entregar. Responde à pergunta ``o que a entrega deve ser e fazer?'', distinta do Escopo do Projeto (o ``como'') e da Qualidade, tratados nas partes seguintes.

\subsection{Descrição do Produto}
\begin{itemize}[itemsep=3pt]
  \item \textbf{O que é:} uma plataforma web de rastreamento e gestão de frotas, acessada por navegador.
  \item \textbf{Para quem:} pequenas transportadoras, operadas por um gestor de frota e utilizadas também pelos motoristas.
  \item \textbf{Como funciona:} integra-se, via API, a rastreadores GPS já existentes nos veículos da transportadora parceira; não há fabricação de hardware próprio. Os dados de posição chegam do provedor, são processados pela plataforma e transformados em mapa, alertas e relatórios.
  \item \textbf{O que entrega:} visibilidade quase em tempo real da frota, alertas automáticos (manutenção, vencimento de documentos, cercas virtuais e perda de sinal) e indicadores de uso para apoiar decisões de gestão baseadas em dados reais.
\end{itemize}

\subsection{Arquitetura e Tecnologias}
Como o produto é uma plataforma de software, ``o que ele é'' inclui as decisões de tecnologia e infraestrutura que o sustentam. As escolhas abaixo são a arquitetura de referência do projeto; mudanças estruturais nela seguem o controle de mudanças definido na Parte II.

\subsubsection*{Pilha de tecnologia}
\begin{center}
\begin{tabularx}{\textwidth}{l L L}
\toprule
\textbf{Camada} & \textbf{Tecnologia} & \textbf{Por que} \\
\midrule
Backend / API & Python com Django e Django REST Framework & Framework maduro, produtivo e adequado a regras de negócio e relatórios; expõe uma API REST consumida pelo frontend. \\
Tarefas agendadas & Job agendado (comando de gerenciamento do Django via cron) & Busca periódica das posições do rastreador e verificação das regras de alerta, sem a complexidade de uma fila de mensagens dedicada. \\
Banco de dados & PostgreSQL & Armazena veículos, motoristas, trajetos e alertas; os cálculos de distância e de cerca virtual (ponto e raio) são resolvidos na própria aplicação. \\
Frontend & React (aplicação de página única) com Leaflet e Recharts & Leaflet renderiza o mapa sobre tiles do OpenStreetMap; Recharts monta os gráficos do painel; a interface é responsiva. \\
Autenticação & Tokens JWT com perfil único (gestor) & Sessões seguras; a primeira versão atende apenas o perfil de gestor. \\
Integração externa & Cliente HTTP para a API do provedor de rastreamento & Consome as posições por chamadas agendadas (polling) ou por webhook, conforme o provedor piloto suportar. \\
Geração de relatórios & Exportação em PDF e XLS & Relatórios entregues em formato imprimível (PDF) e em planilha (XLS) para análise. \\
\bottomrule
\end{tabularx}
\end{center}

\subsubsection*{Infraestrutura}
\begin{itemize}[itemsep=2pt]
  \item \textbf{Contêineres:} a aplicação roda em contêineres Docker orquestrados por docker-compose (serviços: API, banco PostgreSQL e frontend).
  \item \textbf{Borda e segurança de transporte:} proxy reverso Nginx com TLS, servindo tudo sob HTTPS.
  \item \textbf{Hospedagem:} nuvem de baixo custo ou VPS, coerente com a restrição de orçamento zero para hardware próprio.
  \item \textbf{Operação:} segredos em variáveis de ambiente, backup periódico do banco e registro de logs da aplicação.
\end{itemize}

\subsubsection*{Fluxo de dados, ponta a ponta}
\begin{enumerate}[itemsep=2pt]
  \item O rastreador GPS do veículo reporta a posição ao provedor de rastreamento de terceiros.
  \item Um job agendado consome essas posições pela API do provedor e as grava no PostgreSQL.
  \item As regras de alerta são avaliadas sobre os dados recebidos, na mesma execução periódica.
  \item O gestor visualiza mapa, alertas e relatórios no frontend React.
\end{enumerate}

\subsection{Requisitos Funcionais}
Organizados pelos módulos do produto.

\begin{itemize}[itemsep=3pt]
  \item \textbf{Cadastros:} cadastro de veículos (placa, modelo, hodômetro, documentos), de motoristas (dados, CNH e validade) e do vínculo motorista--veículo. O motorista é um registro da frota, não um usuário do sistema nesta fase.
  \item \textbf{Integração e ingestão:} integração com a API do provedor de rastreamento piloto e ingestão periódica das posições recebidas.
  \item \textbf{Monitoramento:} mapa com a localização quase em tempo real dos veículos ativos e histórico de trajetos por veículo e por período.
  \item \textbf{Cercas virtuais:} definição de áreas circulares (ponto central e raio) permitidas ou proibidas por veículo.
  \item \textbf{Motor de alertas:} geração automática dos alertas do catálogo (Seção 1.5), com registro de status (aberto, resolvido) e notificação no aplicativo e por e-mail.
  \item \textbf{Relatórios:} emissão dos relatórios do catálogo (Seção 1.6), com filtros e exportação em PDF e XLS.
  \item \textbf{Painel de indicadores:} veículos ativos e alertas pendentes por categoria.
  \item \textbf{Acesso:} autenticação com perfil de gestor.
\end{itemize}

\subsection{Catálogo de Alertas}
Detalhar ``emitir alertas'' significa dizer exatamente quais alertas, e como cada um é disparado. Em vez de uma lista longa de alertas parecidos, o TrackFleet organiza a emissão em quatro motores de alerta configuráveis, cada um cobrindo vários itens da frota. Cada alerta gerado é registrado na plataforma e notificado ao gestor no aplicativo e por e-mail.

\begin{center}
\begin{tabularx}{\textwidth}{L L L}
\toprule
\textbf{Motor de alerta} & \textbf{Itens que cobre} & \textbf{Gatilho (condição de disparo)} \\
\midrule
Manutenção preventiva (por km/tempo) & Troca de óleo, revisão periódica, pneus e freios & Quilometragem rodada ou tempo decorrido desde a última execução do item atinge o limite configurado \\
Vencimento de documento (por data) & Licenciamento/IPVA, seguro do veículo, CNH do motorista & Data de vencimento do documento se aproxima, com a antecedência configurada \\
Cerca virtual (geofence) & Áreas permitidas ou proibidas por veículo & O veículo cruza o raio de uma área definida (ponto central e raio) \\
Perda de sinal do rastreador & Todos os veículos ativos & Nenhuma posição recebida do veículo por tempo acima do limite \\
\bottomrule
\end{tabularx}
\end{center}

\noindent Os limites (quilometragem, antecedência de vencimento, ponto e raio da cerca, tempo sem sinal) e os itens de cada motor são parâmetros configuráveis pelo gestor, não valores fixos no código. Alertas que dependeriam de telemetria não garantida pela API do provedor (por exemplo, excesso de velocidade e ociosidade de motor) ficam fora desta fase, conforme a Seção 1.7.

\subsection{Catálogo de Relatórios}
Detalhar ``emitir relatórios'' significa dizer quais relatórios e o que cada um responde. Todos são filtráveis por período, veículo e motorista, e exportáveis em PDF e XLS.

\begin{center}
\begin{tabularx}{\textwidth}{L L}
\toprule
\textbf{Relatório} & \textbf{O que responde} \\
\midrule
Relatório de alertas & Quais alertas foram emitidos no período, por categoria, veículo e motorista, com o status de cada um (aberto ou resolvido); inclui os documentos próximos do vencimento \\
Relatório de uso por veículo & Km rodados, horas em operação e número de viagens de cada veículo \\
Relatório de manutenção & Manutenções previstas frente às realizadas por veículo \\
\bottomrule
\end{tabularx}
\end{center}

\noindent O painel gerencial (dashboard), com veículos ativos e alertas pendentes, é uma visão de tela do produto (Seção 1.3), não um relatório exportável.

\subsection{Requisitos Não Funcionais}
\begin{itemize}[itemsep=2pt]
  \item \textbf{Conformidade com a LGPD:} consentimento informado e minimização no tratamento de dados de localização e de motoristas; retenção do histórico por prazo limitado; acesso restrito por autenticação.
  \item \textbf{Atualidade da posição:} a localização exibida no mapa reflete a última leitura do rastreador com atraso perceptível mínimo, respeitado o intervalo de envio do provedor.
  \item \textbf{Segurança:} tráfego sob HTTPS, senhas armazenadas com hash e controle de acesso autenticado.
  \item \textbf{Responsividade:} interface utilizável em desktop e em navegador mobile.
  \item \textbf{Disponibilidade:} plataforma no ar durante toda a janela de testes com a transportadora piloto.
  \item \textbf{Desempenho:} telas principais (mapa e painel) respondem em tempo adequado ao uso do gestor.
\end{itemize}

\subsection{Fora do Escopo do Produto (Exclusões)}
Tão importante quanto listar o que entra é deixar explícito o que não entra, para evitar expectativas equivocadas:
\begin{itemize}[itemsep=2pt]
  \item Fabricação, aquisição ou instalação física de rastreadores GPS próprios.
  \item Aplicativo mobile nativo (iOS/Android); a primeira versão é web responsiva.
  \item Módulo de pagamento, faturamento ou cobrança.
  \item Roteirização ou otimização automática de rotas.
  \item Integração simultânea com múltiplos provedores de rastreamento (apenas um provedor piloto nesta fase).
  \item Bloqueio remoto do veículo ou qualquer atuação sobre o hardware embarcado.
  \item Login e aplicativo próprios para o motorista; nesta fase o motorista é apenas um registro da frota, e o único usuário do sistema é o gestor.
  \item Alertas que dependem de telemetria não garantida pela API do provedor: excesso de velocidade, ociosidade de motor (motor ligado parado) e uso fora do horário de expediente.
  \item Cercas virtuais desenhadas como polígonos livres; nesta fase a cerca é circular (ponto central e raio).
  \item Indicador de economia estimada em reais, por depender de premissas de custo frágeis nesta fase.
\end{itemize}

\subsection{Estrutura do Escopo: Backlog do Produto}
Por se tratar de um projeto conduzido com governança híbrida e leve (conforme o Documento de Governança), o escopo do produto é estruturado como um backlog, decomposto em épicos, features e histórias de usuário, em vez de uma EAP rígida.

\begin{itemize}[itemsep=4pt]
  \item[\textbf{E1}] Monitoramento em tempo real --- localização atual e histórico de trajetos.
  \begin{itemize}
    \item Como gestor, quero ver todos os veículos ativos em um mapa, para saber onde a frota está agora.
    \item Como gestor, quero consultar o trajeto de um veículo em um período, para auditar o uso.
  \end{itemize}
  \item[\textbf{E2}] Alertas automáticos --- manutenção, vencimento de documentos, cercas virtuais e perda de sinal.
  \begin{itemize}
    \item Como gestor, quero receber um alerta antes que a manutenção vire corretiva, para reduzir custo.
    \item Como gestor, quero ser avisado quando um documento estiver perto de vencer, para não rodar irregular.
    \item Como gestor, quero um alerta quando um veículo sair de uma área permitida, para agir na hora.
  \end{itemize}
  \item[\textbf{E3}] Gestão de frota e usuários --- cadastros, vínculos e perfis.
  \begin{itemize}
    \item Como gestor, quero cadastrar veículos e motoristas e vinculá-los, para organizar a frota.
  \end{itemize}
  \item[\textbf{E4}] Relatórios e indicadores --- painel e exportação.
  \begin{itemize}
    \item Como gestor, quero um painel com veículos ativos e alertas pendentes, para decidir rapidamente.
    \item Como gestor, quero exportar o relatório de alertas do período, para prestar contas à diretoria.
  \end{itemize}
\end{itemize}

\subsection{Value Breakdown Structure (VBS)}
A VBS conecta o escopo do produto ao valor pretendido, tornando visível qual entregável move mais o ponteiro do projeto.

\begin{center}
\begin{tabularx}{\textwidth}{L r}
\toprule
\textbf{Entregável (épico)} & \textbf{Valor} \\
\midrule
E1 --- Monitoramento em tempo real & 35\% \\
E2 --- Alertas automáticos & 35\% \\
E4 --- Relatórios e indicadores & 20\% \\
E3 --- Gestão de frota e usuários (infraestrutura base) & 10\% \\
Conformidade com a LGPD & 100\% (obrigatório) \\
\bottomrule
\end{tabularx}
\end{center}

\noindent Leitura da VBS: o monitoramento em tempo real e os alertas automáticos são os entregáveis mais valiosos e devem ser priorizados sobre relatórios avançados caso o prazo do semestre aperte. A conformidade com a LGPD é obrigatória e não é negociável, independentemente de prioridade de valor.

\vspace{2mm}
\hrule
\vspace{4mm}

%======================================================================
\section{Escopo do Projeto (o como)}
%======================================================================

\noindent\textbf{\large Escopo do Projeto}

\vspace{1mm}
\noindent Todo o trabalho necessário, e apenas o necessário, para entregar o Escopo do Produto descrito na Parte I. É o componente mais importante da linha de base, pois nele reside o valor esperado do projeto.

\subsection{Descrição do Trabalho do Projeto}
\begin{itemize}[itemsep=2pt]
  \item Levantar requisitos com a transportadora parceira e mapear a integração com a API do provedor de rastreamento piloto.
  \item Projetar a arquitetura da solução (backend Django, frontend React, banco PostgreSQL, integração externa e infraestrutura em contêineres).
  \item Desenvolver o backend: API própria, ingestão das posições e os quatro motores de alerta do catálogo.
  \item Desenvolver o frontend: mapa, painel de indicadores, cadastros e relatórios exportáveis.
  \item Implementar a autenticação do gestor e as medidas de conformidade com a LGPD.
  \item Testar a solução (testes automatizados e testes com dados reais/simulados do rastreador).
  \item Homologar a plataforma com a transportadora piloto.
  \item Treinar o gestor piloto e coletar feedback.
  \item Documentar o projeto e encerrar a primeira fase, com registro de lições aprendidas.
\end{itemize}

\subsection{Estrutura Analítica do Projeto (EAP)}
A EAP decompõe o trabalho total em pacotes menores, atribuíveis e mensuráveis, tornando o acordo sobre o escopo visível a todos os envolvidos. O projeto está no topo, os grandes entregáveis no segundo nível e os pacotes de trabalho no terceiro.

\begin{center}
\resizebox{\textwidth}{!}{%
\begin{forest}
  for tree={
    draw,
    rounded corners,
    align=center,
    inner sep=5pt,
    l sep=16pt,
    s sep=8pt,
    edge={-},
    font=\small,
  }
  [{1\\TrackFleet}
    [{1.1\\Gestão do\\Projeto}
      [{1.1.1\\Termo de abertura\\e governança}]
      [{1.1.2\\Planejamento e\\acompanhamento}]
    ]
    [{1.2\\Requisitos}
      [{1.2.1\\Entrevistas com a\\transportadora}]
      [{1.2.2\\Especificação de\\requisitos}]
    ]
    [{1.3\\Desenvolvimento}
      [{1.3.1\\Backend, API e\\ingestão}]
      [{1.3.2\\Motor de\\alertas}]
      [{1.3.3\\Frontend (mapa,\\painel, relatórios)}]
      [{1.3.4\\Autenticação (gestor)\\e LGPD}]
    ]
    [{1.4\\Testes e\\Homologação}
      [{1.4.1\\Testes\\automatizados}]
      [{1.4.2\\Piloto com a\\transportadora}]
    ]
    [{1.5\\Treinamento e\\Encerramento}
      [{1.5.1\\Treinamento do\\gestor piloto}]
      [{1.5.2\\Documentação e\\lições aprendidas}]
    ]
  ]
\end{forest}
}
\end{center}

\subsection{Dicionário da EAP}
\begin{center}
\begin{tabularx}{\textwidth}{l L L}
\toprule
\textbf{Pacote} & \textbf{Descrição} & \textbf{Critério de aceite} \\
\midrule
1.3.1 Backend, API e ingestão & API própria em Django consumindo dados do rastreador piloto e persistindo posições no PostgreSQL & Dados do rastreador chegam à plataforma e ficam disponíveis para consulta \\
1.3.2 Motor de alertas & Os quatro motores que avaliam os dados e geram os alertas do catálogo, com status e notificação & Cada motor do catálogo dispara corretamente na condição configurada \\
1.3.3 Frontend & Mapa, painel, cadastros e relatórios exportáveis em React & Gestor localiza um veículo e visualiza alertas pendentes em poucos cliques \\
1.3.4 Autenticação e LGPD & Perfil de gestor e medidas de minimização e proteção de dados & Acesso autenticado e sem exposição indevida de dados pessoais \\
1.4.2 Piloto & Uso real da plataforma pela transportadora parceira por um período de teste & Transportadora usa a plataforma e valida os indicadores gerados \\
\bottomrule
\end{tabularx}
\end{center}

\subsection{Fora do Escopo do Projeto (Exclusões)}
\begin{itemize}[itemsep=2pt]
  \item Compra, instalação física ou manutenção de rastreadores GPS.
  \item Suporte comercial e operação em produção após o encerramento do piloto acadêmico.
  \item Integração com mais de um provedor de hardware de rastreamento.
  \item Certificações formais de segurança ou auditorias externas de compliance além da adequação básica à LGPD.
  \item Desenvolvimento de aplicativo mobile nativo.
\end{itemize}

\subsection{Linha de Base do Escopo e Controle de Mudanças}
A linha de base do escopo deste documento é aprovada pelo patrocinador (Diretoria da transportadora parceira), conforme definido no Documento de Governança. Qualquer mudança de escopo que altere objetivo, prazo ou entregáveis essenciais segue o mesmo fluxo de escalonamento já definido: sem consenso entre os dois GPs, ou se a mudança afetar objetivo/prazo/valor esperado, a decisão sobe ao patrocinador e é registrada no log de decisões do projeto. Isso evita o ``scope creep'', o crescimento não controlado do escopo por pequenos acréscimos aceitos por fora do processo.

\subsection{Indicadores de Escopo}
\begin{center}
\begin{tabularx}{\textwidth}{L l l}
\toprule
\textbf{Indicador} & \textbf{Fórmula} & \textbf{Leitura} \\
\midrule
Precisão da Definição & Itens entregues corretamente $\div$ Itens planejados $\times$ 100 & Maior é melhor \\
Scope Creep & Entregas não planejadas $\div$ Total de entregas $\times$ 100 & Menor é melhor \\
Estabilidade de Requisitos & Requisitos não alterados $\div$ Total de requisitos $\times$ 100 & Maior é melhor \\
\bottomrule
\end{tabularx}
\end{center}

\noindent Esses indicadores devem ser acompanhados a cada check-in do projeto (conforme a cadência definida na Governança), para que desvios de escopo sejam percebidos antes da entrega final, e não durante ela.

\vspace{2mm}
\hrule
\vspace{4mm}

%======================================================================
\section{Qualidade}
%======================================================================

\noindent\textbf{\large Qualidade}

\vspace{1mm}
\noindent A terceira dimensão do escopo: não um domínio à parte, mas um atributo integrante dele. A qualidade foca em atender às expectativas dos stakeholders e em cumprir os requisitos definidos tanto no Escopo do Produto quanto no Escopo do Projeto.

\subsection{Requisitos de Qualidade do Produto}
\begin{itemize}[itemsep=2pt]
  \item Precisão e atualidade da localização exibida no mapa, com latência aceitável em relação à posição real do rastreador.
  \item Confiabilidade dos alertas, com o mínimo possível de falsos positivos e falsos negativos em cada motor do catálogo.
  \item Usabilidade do painel: o gestor deve conseguir localizar um veículo e visualizar alertas pendentes em poucos cliques.
  \item Conformidade com a LGPD no tratamento de dados de localização e de motoristas (consentimento informado e minimização de dados coletados).
  \item Disponibilidade da plataforma durante toda a janela de testes com a transportadora piloto.
\end{itemize}

\subsection{Requisitos de Qualidade do Projeto (Processo)}
\begin{itemize}[itemsep=2pt]
  \item Todo código é revisado pelos dois GPs antes de ser integrado --- não há code review unilateral.
  \item Rotinas críticas (integração com a API do rastreador, ingestão e cálculo de alertas) possuem testes automatizados.
  \item A documentação do projeto é atualizada a cada entrega, não deixada para o final.
\end{itemize}

\subsection{Definição de Pronto (DoD)}
A DoD é o conjunto de critérios que uma funcionalidade precisa cumprir para ser considerada apta ao uso pela transportadora piloto. Ela é definida antes de a funcionalidade começar a ser construída, para blindar o aceite e evitar a conversa em que a equipe diz que terminou e o patrocinador diz que não.

\begin{itemize}[itemsep=3pt]
  \item[\dodbox] Funcionalidade implementada conforme o requisito descrito no Escopo do Produto.
  \item[\dodbox] Código revisado por ambos os GPs (Enzo Allebrand e Kerlon Ribeiro).
  \item[\dodbox] Testada com dados reais ou simulados do rastreador GPS piloto.
  \item[\dodbox] Sem exposição indevida de dados pessoais (checagem de conformidade com a LGPD).
  \item[\dodbox] Validada e aceita pelos dois GPs; mudanças de escopo, prazo ou objetivo passam ainda pelo patrocinador, conforme a matriz RACI da Governança.
\end{itemize}

\subsection{Critérios de Aceite por Entregável}
\begin{center}
\begin{tabularx}{\textwidth}{L L}
\toprule
\textbf{Entregável} & \textbf{Critério de aceite} \\
\midrule
Monitoramento em tempo real (E1) & O gestor visualiza, no mapa, a posição atual de todos os veículos ativos com atraso perceptível mínimo \\
Alertas automáticos (E2) & Cada motor de alerta do catálogo dispara automaticamente na condição configurada, sem falso positivo ou negativo relevante \\
Relatórios e indicadores (E4) & O painel exibe corretamente veículos ativos e alertas pendentes, e os relatórios são exportáveis em PDF e XLS \\
Piloto com a transportadora parceira & A transportadora utiliza a plataforma no período de teste e confirma, por escrito, que os indicadores refletem a realidade da frota \\
\bottomrule
\end{tabularx}
\end{center}

\subsection{Validar \texttimes\ Controlar o Escopo}
As duas atividades são complementares, mas distintas:
\begin{itemize}[itemsep=3pt]
  \item \textbf{Validar o Escopo} --- aceite formal, externo: o patrocinador (Diretoria da transportadora parceira) aprova cada marco relevante do projeto.
  \item \textbf{Controlar o Escopo} --- verificação interna, contínua: os dois GPs conferem, a cada check-in, se o que está sendo construído continua dentro da linha de base aprovada.
\end{itemize}
\noindent Um valida o resultado; o outro cuida do processo. Confundir os dois é um dos erros mais comuns de gestão de escopo.

\subsection{Prevenção de Scope Creep}
As regras de proteção (guardrails) já definidas no Documento de Governança sustentam também a qualidade do escopo:
\begin{itemize}[itemsep=2pt]
  \item Toda mudança relevante de escopo, prazo ou objetivo passa pelo patrocinador antes de ser aceita.
  \item Toda decisão importante é registrada no log de decisões do projeto, no momento em que é tomada.
  \item Nenhuma funcionalidade é considerada ``pronta'' sem passar pela DoD definida na Seção 3.3.
  \item Sem consenso entre os dois GPs, a decisão sobe direto ao patrocinador, que desempata.
\end{itemize}

\subsection{Métricas de Qualidade}
\begin{center}
\begin{tabularx}{\textwidth}{L l}
\toprule
\textbf{Métrica} & \textbf{Leitura} \\
\midrule
Taxa de defeitos encontrados no piloto & Menor é melhor \\
\% de funcionalidades aceitas sem retrabalho na primeira validação & Maior é melhor \\
\% de alertas corretos (sem falso positivo/negativo) & Maior é melhor \\
\bottomrule
\end{tabularx}
\end{center}

\end{document}
